Siedendes Wasser hat die Dichte $\rho_\text{s} = \rhWsO$ (FoTaBe S.~192), die Verdampfungswärme ist $L_\text{v} = \LvWO$ (FoTaBe S.~191):
\begin{align}
 Q &= \QvF = \rhWs\cdot\VWs\cdot\LvW = \Qv \approx \result{\QvP{6}{2}}
\end{align}
Mit dieser Wärme kann man die Wassermasse
\begin{align}
 m' &= \mWaF = \frac{\Qv}{\cW\cdot\left(\TsO-\TnO\right)} = \mWa
\end{align}
von $\vartheta_0 = \TnO$ auf $\vartheta_\text{s} = \TsO$ erwärmen, mit $\rho_0 = \rhWnO$ (FoTaBe S.~192) also das Volumen
\begin{align}
 V' &= \VWnF = \frac{\mWa}{\rhWn} = \VWnl \approx \result{\VWnlP{0}{2}}
\end{align}
Das Verdampfen braucht gut fünfmal so viel Wärme wie das Erwärmen derselben Wassermenge von \TnO\ bis zum Sieden.
